\section{Round-twenty-three reconstruction: local phase Laplace theory and shared-policy labelled contractions}
\label{sec:r23-d1}

An LDP determines exponential phase costs but not polynomial phase weights.
The active synthesis therefore begins with a local density/coefficient theorem
near every phase minimum and applies a Morse--Bott Laplace expansion.  Phase
cells are chosen as rate-continuity sets, and exponential tightness supplies
the missing separation from their noncompact complement.  Control is then
performed on the full phase-posterior state with one policy shared by all
components.

\subsection{Microscopic phase events and local normal form}

Let $\Omega_N$ be the original microscopic sample space and let
$M_N:\Omega_N\to\mathbb R^d$ be a continuous finite-dimensional order
parameter built from the A or B observables.  The upstream source-inserted A2
or B1 local theorem gives, in a neighbourhood of every regular minimizer, the
joint local form
\[
 P_N\{M_N\in dm,Z_N\in dz\}
 =N^{-\lambda_j}e^{-N J_j(m,z)}
 \bigl(a_j(m,z)+O(N^{-1/2})\bigr)\,dm\,\varkappa_N(dz),   \tag{D1.1}
\]
where $\varkappa_N$ is the appropriate lattice/continuous reference, $a_j$ is
positive and $C^2$, and the error is uniform with two derivatives.  This is a
strictly stronger input than an LDP and is proved by the same source-inserted
Fourier inversion used upstream.

Assume the minima of $m\mapsto J_j(m,z_j)$ form finitely many disjoint compact
$C^3$ manifolds $\mathcal M_j$ of dimension $r_j$.  In a tubular chart
$m=(y,n)$, $y\in\mathcal M_j$ and $n$ normal,
\[
 J_j(y,n,z_j)=\gamma_j+rac12
 \langle H_j(y)n,n\rangle+O(|n|^3),\qquad H_j(y)>cI.      \tag{D1.2}
\]
Choose disjoint tubular neighbourhoods $U_j$ whose boundaries are regular
level sets of the squared normal distance and satisfy
\[
 \inf_{\partial U_j}J\ge\gamma_j+3\eta.                  \tag{D1.3}
\]
The exact phase event is
\[
 \mathcal C_{j,N}=\{\omega\in\Omega_N:M_N(\omega)\in U_j\}.             \tag{D1.4}
\]
It is a Borel event in the original microscopic law; no latent phase is added
before this disintegration.

\begin{lemma}[Separation including the noncompact complement]
\label{lem:r23-d1-separation}
There is a compact set $K$ and $\eta>0$ such that
\[
 \inf_{m\notin K}J(m)\ge\gamma_*+4\eta,
 \qquad
 \inf_{m\in K\setminus\cup_jU_j}J(m)\ge\gamma_*+2\eta, \tag{D1.5}
\]
where $\gamma_*=\min_j\gamma_j$.  Each $U_j$ is an
$I$-continuity set for the phase-conditioned LDP.
\end{lemma}

\begin{proof}
Goodness of the upstream rate and exponential tightness allow a compact
sublevel $K$ containing every minimizer with the first bound.  Inside $K$, the
minimizer manifolds are compact and disjoint.  Shrink their tubular
neighbourhoods until (D1.3) holds; the remaining compact set contains no
minimizer, so its rate is separated by a positive amount.  The boundary rate
is strictly larger than the phase minimum, hence the infimum over the closure
and interior of $U_j$ agree.  This is the required conditioning regularity.
\end{proof}

\subsection{Morse--Bott phase weights}

\begin{theorem}[Phase weight expansion]
\label{thm:r23-d1-weights}
For every phase cell (D1.4),
\[
 P_N(\mathcal C_{j,N})
 =e^{-N\gamma_j}N^{-\kappa_j}
 \bigl(c_j+O(N^{-1/2})\bigr),                             \tag{D1.6}
\]
where
\[
 \kappa_j=\lambda_j+{d-r_j\over2},
 \qquad
 c_j=(2\pi)^{(d-r_j)/2}
 \int_{\mathcal M_j}{a_j(y,0,z_j)\over
       \sqrt{\det H_j(y)}}\,d\operatorname{vol}_{\mathcal M_j}(y)>0.  \tag{D1.7}
\]
All lattice, exact-number, saddle, and conditioning powers are included in
$\lambda_j$ from the upstream local coefficient; none is inferred from the
LDP.
\end{theorem}

\begin{proof}
Use tubular coordinates in (D1.1), set $n=N^{-1/2}\xi$, and apply (D1.2).
The Jacobian is $1+O(N^{-1/2}|\xi|)$ and the cubic remainder is uniformly
integrable against the Gaussian on $|\xi|\le N^{1/10}$.  Dominated convergence
gives the normal Gaussian determinant and the integral over
$\mathcal M_j$.  The complement of that tube is exponentially smaller by
(D1.3), and the complement of all phase cells is exponentially smaller by
Lemma~\ref{lem:r23-d1-separation}.  This is the standard Morse--Bott Laplace
calculation \cite{Wong2001}, now applied to the actual local coefficient
(D1.1).
\end{proof}

The exact disintegration
\[
 P_N=\sum_jw_{j,N}P_{j,N}+P_N(\,\cdot\,;\,M_N\notin\cup_jU_j),
 \qquad
 w_{j,N}=P_N(\mathcal C_{j,N}),                           \tag{D1.8}
\]
has an exponentially negligible remainder and the explicit weights (D1.6).

\subsection{Labelled and conditional large deviations}

Let $J_N=j$ on $\mathcal C_{j,N}$ and $J_N=\dagger$ on the remainder.  By the
$I$-continuity from Lemma~\ref{lem:r23-d1-separation}, the conditional laws
$P_{j,N}$ satisfy the component LDP with rate
\[
 I_j(z)=I(z,j)-\gamma_j.                                 \tag{D1.9}
\]
The labelled pair $(J_N,Z_N)$ has good rate
$\gamma_j+I_j(z)-\gamma_*$, and contraction over $j$ gives
\[
 I(z)=\min_j\{\gamma_j-\gamma_*+I_j(z)\}.                \tag{D1.10}
\]
The polynomial weights in (D1.6) affect coexistence mixtures but not the
exponential rate (D1.10); the two statements are kept distinct.

\subsection{Phase-conditioned belief state and one shared policy}

For each phase let $\pi_j$ be its C1 hidden-state belief and let
$w=(w_j)$ be the phase posterior.  The sufficient state is
\[
 \mathbf b=(w_1,\pi_1;\ldots;w_J,\pi_J).                 \tag{D1.11}
\]
Given one common action $a$, phase $j$ has observation evidence
$r_j(\pi_j,a,y)$ and posterior $B_j(\pi_j,a,y)$.  The common evidence and
Bayes update are
\[
 r(\mathbf b,a,y)=\sum_jw_jr_j(\pi_j,a,y),
 \quad
 w_j'={w_jr_j\over r},
 \quad
 \pi_j'=B_j(\pi_j,a,y),                                  \tag{D1.12}
\]
with an arbitrary convention on the zero-evidence set, which has zero weight.
C1's integrated Feller theorem proves continuity of the kernel on the
augmented state.

For a fixed admissible policy $\alpha$ let
\[
 \mathcal V_N^\alpha(\mathbf b)
 =\log\sum_jw_jE_{j}^{\alpha}
       \exp\{N F(Z_N)\}.                                 \tag{D1.13}
\]
The controlled value is
\[
 \mathcal V_N(\mathbf b)=\sup_{\alpha\text{ common}}
 \mathcal V_N^\alpha(\mathbf b).                         \tag{D1.14}
\]
Optimization is outside the phase sum.  Applying separate optimizers inside
(D1.13) would give a controller that observes the unobserved phase and is not
used.

\begin{theorem}[Shared-policy Bellman principle]
\label{thm:r23-d1-control}
The augmented state (D1.11)--(D1.12) is Markov and its exponential value
satisfies the Bellman recursion with one action common to every phase.  At the
large-deviation scale,
\[
 \lim_{N\to\infty}{1\over N}\mathcal V_N(\mathbf b)
 =\sup_{\alpha\text{ common}}
 \max_j\left\{-\gamma_j+
     \sup_z(F^\alpha_j(z)-I_j^\alpha(z))
ight\}.         \tag{D1.15}
\]
Finite-$N$ dynamic consistency is retained by the full posterior state, not by
the marginal log-sum alone.
\end{theorem}

\begin{proof}
Condition on the next common observation.  Formula (D1.12) is Bayes' rule for
the exact mixture and gives the tower property for (D1.13).  Compact actions
and the integrated Feller kernel yield measurable maximizers.  For fixed
$\alpha$, apply the component Laplace principles and the elementary finite
log-sum asymptotic.  Only then take the supremum over common policies.  This
proves (D1.15) without interchanging phasewise optimization and aggregation.
\end{proof}

\subsection{Complex phase charts and genuine poles}

The source-inserted local theorem gives, for complex $z$ near zero,
\[
 Z_{j,N}(z)=e^{N\Lambda_j(z)}N^{-\kappa_j}
 a_j(z)(1+O(N^{-1/2})),                                  \tag{D1.16}
\]
with $a_j(0)=c_j>0$ and analytic $a_j$.  Shrinking the chart makes $a_j$
zero-free.  In a region where one phase $j_*$ has
\[
 \operatorname{Re}\Lambda_{j_*}(z)
 \ge\max_{k\ne j_*}\operatorname{Re}\Lambda_k(z)+g,   \tag{D1.17}
\]
Rouche's theorem shows that the full sum is zero-free for large $N$ and its
logarithmic derivatives agree with phase $j_*$ up to $O(e^{-Ng})$.  At a
coexistence surface, (D1.17) is not asserted; possible zeros are described by
the finite analytic sum (D1.16).  Thus complex dominance is derived from the
upstream expansion and is not assumed globally.

\subsection{Labelled Gaussian coexistence}

Let $m_j$ and $\Sigma_j$ be the component mean and covariance obtained from the
upstream local theorem.  Under $P_{j,N}$,
\[
 \sqrt N\,(X_N-m_j)\Rightarrow N(0,\Sigma_j).             \tag{D1.18}
\]
For the labelled variable
$\bigl(J_N,\sqrt N(X_N-m_{J_N})\bigr)$, the limiting mixture weights are the
normalised versions of
\[
 e^{-N\gamma_j}N^{-\kappa_j}c_j.                         \tag{D1.19}
\]
At an exponential tie, phases with smaller $\kappa_j$ dominate; among phases
with equal $\kappa_j$, the weights are proportional to $c_j$.  Every local
partition, saddle, and conditioning prefactor is already present in
$\kappa_j,c_j$ through Theorem~\ref{thm:r23-d1-weights}.  The unlabelled
variable is obtained only after this phasewise centering; it is not incorrectly
assigned the phase means after a second centering.

\begin{theorem}[Coexistence mixture]
\label{thm:r23-d1-gaussian}
Along any subsequence on which the normalised weights (D1.19) converge to
$(\rho_j)$,
\[
 \mathcal L\bigl(J_N,\sqrt N(X_N-m_{J_N})\bigr)
 \Rightarrow\sum_j\rho_j\,\delta_j\otimes N(0,\Sigma_j).\tag{D1.20}
\]
The same weights enter the phase-posterior initialization of
Theorem~\ref{thm:r23-d1-control}.
\end{theorem}

\begin{proof}
Apply the component local central limit theorem inside each $I$-continuity
cell and multiply by the phase probability (D1.6).  The complement is
exponentially negligible by Lemma~\ref{lem:r23-d1-separation}.  Normalization
gives (D1.20).
\end{proof}

\subsection{Round-Twenty-Two regression}

Polynomial weights are obtained from the local normal form (D1.1)--(D1.2), not
from an LDP.  Phase boundaries are rate-separated and the noncompact remainder
is controlled by exponential tightness.  Conditional component principles use
$I$-continuity.  One policy is fixed before aggregation.  Complex zero-free
charts require the proved dominance gap (D1.17).  The Gaussian variable is
centered by its own labelled phase and the tie exponent contains every local
prefactor.
